\section{Conclusion}
\label{sec:conclusion}

We tested whether a learned, versioned lifecycle vocabulary helps object-level memory under changes made outside the robot's view, holding the front end, recall, executor, training recipe and evaluation fixed and comparing against the same learned association without that vocabulary.
Under a protocol frozen before validation and test, with a single test read, the answer for this training procedure (five seeds) is yes for near-ideal instance masks on both stale entities and identity continuity, and yes for stale entities only with SAM~2.1 masks.
In a descriptive ablation, identity continuity depends on keeping the versions of retracted entities: deleting them leaves slightly fewer stale entities but loses the identities.
Against the rule-based adapters there is no single winner: with instance masks the persistence-filter adapter leaves fewer stale entities, while most of the in-scope entities it retracts belong to objects that are still in place.
The low absolute identity continuity, especially with a learned segmenter, is the main open problem; the released ledger, exports and index are meant to make that problem measurable.
